\documentclass[no-page-number]{nanzan-abstract}

\usepackage{url}
\usepackage{graphicx}
\usepackage{tikz}
\usepackage{array}
\usepackage{booktabs}
\usetikzlibrary{arrows.meta,positioning}

\setlength{\tabcolsep}{2.5pt}

\ModifyHeading{section}{
  font={\large\bfseries},
  format={\vbox{\noindent #1#2}},
  lines=2,
  before_space=0.8ex,
  after_space=0.5ex
}

\AtBeginDocument{\fontsize{9.5pt}{11.3pt}\selectfont}

\title{姿勢推定を用いた筋力トレーニング評価システムにおける処理の共通化}
\author{2023TS030 川瀬 翔大}
\professor{宮澤 元}

\begin{document}
\maketitle

\section{はじめに}

筋力トレーニングでは，フォームが崩れると狙った部位に十分な負荷をかけにくくなり，トレーニングの効率が低下する可能性がある．しかし，運動中のフォームを客観的に確認することは難しい．そこで，カメラ映像から人体の姿勢を推定し，関節の位置や角度を利用して運動フォームを評価する研究が行われている．このようなシステムでは，ランドマークから角度などを計算し，設定した条件と比較してフォームを判定できる．

一方で，対応する筋力トレーニングの種目を増やすと，種目ごとに評価する関節や目標角度，許容範囲，比較条件などが変わる．そのため，種目を追加するたびに角度計算や判定に関するプログラムを修正する必要があり，対応する種目が増えるほど修正や動作確認の手間も増加する．例えば，スクワットでは股関節・膝・足首を用いて膝関節の角度を評価するのに対し，アームカールでは肩・肘・手首を用いて肘関節の角度を評価する．対象となる関節は異なるが，3点から角度を求め，基準値と比較して判定するという処理には共通する部分がある．

そこで本研究では，種目追加時に既存プログラムを修正する負担を抑えつつ，評価処理を拡張できる構成を明らかにすることを目的とする．具体的には，種目ごとに異なる評価条件を設定データにまとめ，姿勢推定後の角度計算や判定など，複数の種目で利用できる処理を共通化する．これにより，種目固有の情報を設定として分離し，新しい種目へ対応できる構成を検討する．RQは「関節指定，目標角度，判定条件などの評価規則を，どこまで設定データの変更だけで扱えるか．また，設定だけでは扱えない規則にはどのような種類があり，どの段階で共通コードの追加・変更や個別処理が必要になるか」である．

\section{研究の背景}

田邉らは，姿勢推定を用いてサイドレイズやアームカールの運動を3段階に分け，それぞれに応じたフィードバックを行っている[1]．本研究では，運動中の動きを評価する方法の例として参考にする．特に，運動を単に一つの姿勢として評価するのではなく，運動の進み方を考えて評価している点を参考にする．

Mercadal-Baudartらは，単眼カメラによる3次元姿勢推定を用いて複数の運動を計測し，関節角度を評価に利用している[2]．本研究では，姿勢推定によって得られた身体の特徴点から関節の角度を計算し，運動の評価に利用する方法を参考にする．

また，Patidarらは，MediaPipeとOpenCVを用いて，スクワット，アームカール，デッドリフト，サイドレイズなどを関節角度から評価している[3]．本研究では，複数の筋力トレーニング種目に対して関節角度を利用した評価を行う例として参考にする．

これらの研究を参考に，本研究ではフォームの評価精度そのものではなく，複数の種目に対応するときのプログラム構成に着目する．姿勢推定から得られたランドマークを使う処理と種目ごとの評価条件を分離し，同じ処理をどこまで利用できるかを調べる．

\section{研究概要と設計}

本研究では，スクワットとアームカールを対象に，同一の評価処理を適用できる構成を検討する．スクワットでは股関節・膝・足首の3点から膝関節の角度を求め，アームカールでは肩・肘・手首の3点から肘関節の角度を求める．この2種目では関節やランドマークは異なるが，3点から角度を求めて基準値と比較する処理は共通している．図1に想定するシステム構成を示す．

\begin{figure}[t]
\centering
\begin{tikzpicture}[
  node distance=5mm and 8mm,
  box/.style={draw, rounded corners=2pt, align=center, minimum width=42mm, minimum height=9mm, font=\small},
  cfg/.style={draw, rounded corners=2pt, align=left, minimum width=45mm, minimum height=18mm, font=\small},
  arr/.style={-{Latex[length=2mm]}, thick}
]
\node[box] (cam) {カメラ映像};
\node[box, below=of cam] (pose) {姿勢推定\\ランドマーク};
\node[box, below=of pose] (points) {共通処理\\設定から評価対象の3点を取得};
\node[box, below=of points] (feat) {共通処理\\3点から特徴量（角度）を計算};
\node[box, below=of feat] (judge) {共通処理\\基準値・許容範囲・比較条件で判定};
\node[box, below=of judge] (out) {評価結果\\表示・フィードバック};
\node[cfg, right=of feat] (config) {種目別設定\\関節指定\\基準値・許容範囲\\比較条件};
\draw[arr] (cam) -- (pose);
\draw[arr] (pose) -- (points);
\draw[arr] (points) -- (feat);
\draw[arr] (feat) -- (judge);
\draw[arr] (judge) -- (out);
\draw[arr] (config.west) -| (points.east);
\draw[arr] (config.west) -| (feat.east);
\draw[arr] (config.west) -| (judge.east);
\end{tikzpicture}
\caption{想定するシステム構成}
\label{fig:system}
\end{figure}

カメラ映像から姿勢推定を行い，ランドマークを取得する．MediaPipeから取得したランドマークの画像平面上の2次元座標を用いて関節角度を計算する．現在の実装では，MediaPipeの正規化されたx，y座標をそのまま用いて角度を計算しているため，今後は画像幅・高さを用いて画像平面上の座標へ変換してから角度計算する．その後，種目ごとの設定を使って評価に使用する3点を選び，角度を目標角度，許容範囲，比較条件と比べてフォームを判定する．このように，姿勢推定によるランドマーク取得以降の処理を共通化し，種目によって変化する情報を設定として与える構成を想定する．

設定データには，(a)角度計算に使う3点のランドマーク名，(b)目標角度，(c)許容範囲，(d)比較条件を記述する．例えばスクワットでは股関節・膝・足首を指定し，膝角度90度，許容範囲$\pm15$度のように設定する．アームカールでは3点を肩・肘・手首に変更することで，同じ角度計算や判定処理を使えるようにする．種目追加時に設定データの変更だけで対応できる範囲を確認する．

本研究で重要なのはJSONという形式そのものではなく，設定だけで対応できる部分と，プログラムを追加する必要がある部分を分けることである．複数の角度や動作の順番など，現在の設定だけでは対応できない評価方法については，共通処理または個別処理として扱う．この設定，共通処理，個別処理の境界を明確にすることが，本研究の検討対象となる．

\section{実装}

実装にはPython，OpenCV，MediaPipeを用いる．OpenCVはカメラ映像の取得や画像処理，MediaPipe Poseは人体のランドマークの取得に利用する[4,5]．現在はスクワットを対象とした試作を行っており，Webカメラ映像の取得，姿勢推定，ランドマーク・骨格の表示，3点からの膝角度の計算，基準値との比較・判定，結果の表示，JSONからの設定値の読み込みまでを確認している．実装では左右それぞれについて股関節・膝・足首のランドマークを取得し，3点のx，y座標から膝角度を計算している．また，JSONから左右の目標角度と許容誤差を読み込んで判定に利用している．

現在の実装では，角度計算にMediaPipeから得られるランドマークの正規化されたx，y座標を使用している．そのため，今後は画像幅と画像高さを用いて画像平面上の2次元座標へ変換してから角度を計算する．

ただし，現時点では共通化が完成したわけではない．現在のコードでは角度計算に使う股関節・膝・足首がコード内に固定されており，設定を変更するだけで対象関節を切り替えることはまだできていない．今後は，角度計算に使う3点を設定ファイルから読み込み，スクワットとアームカールで同じ処理を使えるようにする．さらに，種目ごとに異なる条件を設定として記述できるようにし，設定の変更だけで対応できる範囲を確認する．

\section{検証方法}

検証では，まず3章で示した設定形式と共通処理を評価対象として固定する．スクワットとアームカールの設定を用意し，同じプログラムで両方の種目を扱えるか確認する．また，計算した角度が正しいか，設定した条件に対して正しく判定できるかを確認する．

各評価方法を追加した際に，設定の変更だけで対応できるか，共通処理の変更が必要か，種目固有の処理が必要かを記録する．対象とするのは，(i)関節の指定だけが違う場合，(ii)複数の角度や条件を組み合わせる場合，(iii)角度の変化や動作の順番など，時間による変化を使う場合である．

\begin{table}[t]
\centering
\caption{検証する評価方法と確認内容}
\small
\begin{tabular}{p{0.25\columnwidth}p{0.25\columnwidth}p{0.40\columnwidth}}
\toprule
評価方法 & 主な違い & 確認する内容\\
\midrule
関節角度 & 対象3点 & 設定だけで変更できるか\\
複数条件 & 複数角度 & 共通処理の変更で対応できるか\\
時間的変化 & フレーム間の関係 & 個別処理が必要か\\
\bottomrule
\end{tabular}
\end{table}

例えば，スクワットの膝角度とアームカールの肘角度では使用する3点は異なるが，「3点から角度を計算して基準値と比べる」という処理は同じである．この場合に設定を変えるだけで動作すれば，関節の指定を設定に分ける方法を確認できる．一方，複数の角度を同時に判定する場合は複数角度を扱う処理が必要になる．また，動作の順番を判定する場合は複数フレームの情報を使う必要がある．これらについて，設定，共通処理，個別処理のどの段階で対応するかを記録する．

角度計算と判定の正しさについても確認する．例えば既知の3点の座標を入力し，90度など期待する角度が得られることを確認する．さらに，目標90度，許容範囲$\pm15$度と設定し，75度，90度，105度などの値を用いて判定結果が設定どおりになるかを確認する．なお，角度が設定した範囲に入ったことは，あくまで設定した基準に対する判定であり，それだけで「正しいフォーム」と判断するものではない．

さらに，種目追加時に変更した部分と追加したコード量を記録し，共通化によってプログラムの修正をどの程度減らせるかを確認する．

\section{考察の観点}

検証結果から，設定を変更するだけで対応できた評価方法と，プログラムの変更が必要だった評価方法を整理する．設定だけでは対応できなかった場合には，新たに必要となった処理を共通処理として追加するべきか，種目固有の処理とするべきかを検討する．例えば，複数の角度を設定できるようにすることで共通処理として対応できる場合がある一方，動作の順番など複雑な処理まで設定に入れると設定自体が分かりにくくなる可能性がある．

そのため，設定に入れる部分とプログラムとして分ける部分をどのように決めるかを考察する．また，設定項目を増やすことで対応できる評価方法が増える一方で，設定が複雑になる可能性もあるため，どこまでを設定で扱うかを検討する．さらに，種目追加時に実際にどの部分を変更したかを確認し，共通化によって既存プログラムへの変更をどの程度減らせたかを確認する．

また，設定だけでは対応できなかった処理について，共通処理と個別処理のどちらにするべきかを整理し，設定の柔軟性と複雑さのバランスを考察する．

\section{おわりに}

本研究では，カメラ映像から姿勢を推定してフォームを評価するシステムを複数の種目に対応させる際の，プログラムの修正や保守の負担に着目した．種目ごとに異なる部分を設定にまとめ，それ以外の処理を共通化する構成を考えた．現在はスクワットの試作を行い，映像取得，姿勢推定，角度計算，判定，結果表示，JSONの読み込みまでを確認している．今後は，画像平面上の座標を用いた角度計算への整理と，角度計算に使う関節を設定から変更できる共通化を行う．その上で，スクワットとアームカールを同じ処理で扱えるかを確認する．

さらに，複数の角度や条件，時間による動作の変化など，評価方法を段階的に変えながら検証し，設定だけで対応できる範囲と，プログラムの変更が必要になる範囲を明らかにする．また，設定だけでは対応できなかった処理について，共通処理として追加する場合と種目固有の処理とする場合を整理する．これらの結果から，種目追加時の変更量を確認し，今回考えた構成によって既存プログラムの修正や保守の負担をどの程度減らせるかを調べる．

\begin{thebibliography}{9}
\small
\setlength{\itemsep}{0pt}

\bibitem{tanabe}
田邉英介 他：
運動過程をフィードバックする筋力トレーニング支援システム，
2021年度電気・情報関係学会九州支部連合大会講演論文集，
p.228，2021．

\bibitem{merca}
C. Mercadal-Baudart, et al.:
Exercise quantification from single camera view markerless 3D pose estimation,
Heliyon, Vol.10, No.6, e27596, 2024.

\bibitem{patidar}
K. Patidar, et al.:
Framework for Gym Posture Estimation Using MediaPipe and OpenCV,
2025 IEEE International Conference on Recent Advances in Computing and Systems (REACS),
DOI: 10.1109/REACS67479.2025.11413297, 2025.

\bibitem{mediapipe}
Google AI Edge: MediaPipe Pose Landmarker,
\url{https://ai.google.dev/edge/mediapipe/solutions/vision/pose_landmarker}
（参照 2026-09-21）．

\bibitem{opencv}
OpenCV: OpenCV Documentation,
\url{https://docs.opencv.org/4.x/}
（参照 2026-09-21）．

\end{thebibliography}

\end{document}
